\documentclass[10pt,journal,compsoc]{IEEEtran}

% --- 基础宏包引用 ---
\usepackage[utf8]{inputenc}   % 原文写作 \usepackage[utf8]{utf8}，utf8 不是宏包名，会编译失败
\usepackage[T1]{fontenc}
\usepackage{amsmath,amssymb,amsfonts}
\usepackage{booktabs}
\usepackage{cite}
\usepackage{microtype}
\usepackage{url}
\usepackage[hidelinks]{hyperref}

% --- 自动生成与管理 BibTeX 参考文献文件 ---
\usepackage{filecontents}
\begin{filecontents*}{\jobname.bib}
@article{xie2018temperature,
  title={Temperature scenario generation for probabilistic load forecasting},
  author={Xie, Jingrui and Hong, Tao},
  journal={IEEE Transactions on Smart Grid},
  volume={9}, number={3}, pages={1680--1687}, year={2018}, publisher={IEEE}
}

@article{liu2017probabilistic,
  title={Probabilistic load forecasting via quantile regression averaging on sister forecasts},
  author={Liu, Bidong and Nowotarski, Jakub and Hong, Tao and Weron, Rafa{\l}},
  journal={IEEE Transactions on Smart Grid},
  volume={8}, number={2}, pages={730--737}, year={2017}, publisher={IEEE}
}

@phdthesis{hong2010short,
  title={Short term electric load forecasting},
  author={Hong, Tao},
  year={2010}, school={North Carolina State University}
}

@article{hong2016probabilistic,
  title={Probabilistic electric load forecasting: A tutorial review},
  author={Hong, Tao and Fan, Shu},
  journal={International Journal of Forecasting},
  volume={32}, number={3}, pages={914--938}, year={2016}, publisher={Elsevier}
}

@article{li2020combining,
  title={Combining probability density forecasts for power electrical loads},
  author={Li, Tianyu and Wang, Yi and Zhang, Ning},
  journal={IEEE Transactions on Smart Grid},
  volume={11}, number={2}, pages={1679--1690}, year={2020}, publisher={IEEE}
}

@article{zhang2019improved,
  title={An improved quantile regression neural network for probabilistic load forecasting},
  author={Zhang, Wenjie and Quan, Hao and Srinivasan, Dipti},
  journal={IEEE Transactions on Smart Grid},
  volume={10}, number={4}, pages={4425--4434}, year={2019}, publisher={IEEE}
}

@article{cannon2011quantile,
  title={Quantile regression neural networks: Implementation in {R} and application to
         precipitation downscaling},
  author={Cannon, Alex J.},
  journal={Computers \& Geosciences},
  volume={37}, number={9}, pages={1277--1284}, year={2011}, publisher={Elsevier}
}

@article{koenker1978regression,
  title={Regression quantiles},
  author={Koenker, Roger and Bassett, Gilbert},
  journal={Econometrica},
  volume={46}, number={1}, pages={33--50}, year={1978}
}

@article{yang2020bayesian,
  title={Bayesian deep learning-based probabilistic load forecasting in smart grids},
  author={Yang, Yandong and Li, Wei and Gulliver, T. Aaron and Li, Shufang},
  journal={IEEE Transactions on Industrial Informatics},
  volume={16}, number={7}, pages={4703--4713}, year={2020}, publisher={IEEE}
}

@article{brusaferri2022probabilistic,
  title={Probabilistic electric load forecasting through {B}ayesian mixture density networks},
  author={Brusaferri, Alessandro and Matteucci, Matteo and Spinelli, Stefano and Vitali, Andrea},
  journal={Applied Energy},
  volume={309}, pages={118341}, year={2022}, publisher={Elsevier}
}

@article{faustine2022fpseq2q,
  title={{FPSeq2Q}: Fully parameterized sequence to quantile regression for net-load
         forecasting with uncertainty estimates},
  author={Faustine, Anthony and Pereira, Lucas},
  journal={IEEE Transactions on Smart Grid},
  volume={13}, number={3}, pages={2440--2451}, year={2022}, publisher={IEEE}
}

@article{li2023probabilistic,
  title={Probabilistic forecasting method for mid-term hourly load time series based on an
         improved temporal fusion transformer model},
  author={Li, Dan and Tan, Ya and Zhang, Yuanhang and Miao, Shuwei and He, Shuai},
  journal={International Journal of Electrical Power \& Energy Systems},
  volume={146}, pages={108743}, year={2023}, publisher={Elsevier}
}

@article{lim2021temporal,
  title={Temporal fusion transformers for interpretable multi-horizon time series forecasting},
  author={Lim, Bryan and Ar{\i}k, Sercan {\"O}. and Loeff, Nicolas and Pfister, Tomas},
  journal={International Journal of Forecasting},
  volume={37}, number={4}, pages={1748--1764}, year={2021}, publisher={Elsevier}
}

@article{wang2019combining,
  title={Combining probabilistic load forecasts},
  author={Wang, Yi and Zhang, Ning and Tan, Yuxiao and Hong, Tao and Kirschen, Daniel S.
          and Kang, Chongqing},
  journal={IEEE Transactions on Smart Grid},
  volume={10}, number={4}, pages={3664--3674}, year={2019}, publisher={IEEE}
}

@article{hong2016gefcom,
  title={Probabilistic energy forecasting: {G}lobal energy forecasting competition 2014
         and beyond},
  author={Hong, Tao and Pinson, Pierre and Fan, Shu and Zareipour, Hamidreza and
          Troccoli, Alberto and Hyndman, Rob J.},
  journal={International Journal of Forecasting},
  volume={32}, number={3}, pages={896--913}, year={2016}, publisher={Elsevier}
}

@article{gneiting2007strictly,
  title={Strictly proper scoring rules, prediction, and estimation},
  author={Gneiting, Tilmann and Raftery, Adrian E.},
  journal={Journal of the American Statistical Association},
  volume={102}, number={477}, pages={359--378}, year={2007}
}

@article{winkler1972decision,
  title={A decision-theoretic approach to interval estimation},
  author={Winkler, Robert L.},
  journal={Journal of the American Statistical Association},
  volume={67}, number={337}, pages={187--191}, year={1972}
}

@article{nowotarski2018recent,
  title={Recent advances in electricity price forecasting: A review of probabilistic
         forecasting},
  author={Nowotarski, Jakub and Weron, Rafa{\l}},
  journal={Renewable and Sustainable Energy Reviews},
  volume={81}, pages={1548--1568}, year={2018}, publisher={Elsevier}
}

@article{baringhaus2004new,
  title={On a new multivariate two-sample test},
  author={Baringhaus, Ludwig and Franz, Carsten},
  journal={Journal of Multivariate Analysis},
  volume={88}, number={1}, pages={190--206}, year={2004}, publisher={Elsevier}
}
\end{filecontents*}

\begin{document}

\title{Probabilistic Electric Load Forecasting Under Limited Historical Samples}
\author{Author Name}

\maketitle

\section{Related Work}

\subsection{Probabilistic Electric Load Forecasting Baselines}
Hong and Fan \cite{hong2016probabilistic} organize probabilistic electric load forecasting
(PLF) methodologies into three sources of uncertainty: simulating the predictors, using a
probabilistic model directly, and converting or combining point forecasts. Four methods are
reviewed below, one representing each of the principal paradigms that have developed within
this taxonomy: temperature scenario generation, quantile regression averaging over a family
of point forecasts, combination of predictive densities, and direct estimation of the
conditional quantile function by a neural network.

Throughout, let $y_t$ denote the realized load at time $t$ and $\hat{y}_{t,q}$ the forecast of
its $q$-th quantile.

\subsubsection{Temperature Scenario-Based Probabilistic Forecasting}
Scenario-based methods quantify load uncertainty by propagating meteorological variability
through a deterministic forecasting engine. Xie and Hong \cite{xie2018temperature} evaluate
three practical strategies for generating temperature scenarios. The \emph{fixed-date} method
assigns the temperature profile of a past year to the target calendar date by date, producing
$k$ scenarios from $k$ years of history. The \emph{shifted-date} method extends this by
additionally sliding each historical profile forward and backward by $1,\dots,n$ days,
yielding $(2n+1)k$ scenarios; because each scenario remains an observed temperature
trajectory that has only been relabelled, the autocorrelation structure of the original
series is retained. The \emph{block bootstrap} method cuts each historical year into $m$-day
blocks and draws each block of the target calendar independently from any of the $k$ years.
The study reports the shifted-date method to give the lowest quantile score, with the
bootstrap not improving on fixed-date, and offers the rule of thumb
$n = \lceil 30/k \rceil - 1$ for selecting the shift range.

Each scenario is passed through a deterministic multiple linear regression model---Tao's
Vanilla benchmark \cite{hong2010short}, the Hong-2014 model, or a reduced T-cube
specification---and the resulting ensemble of point forecasts is converted to percentiles by
the empirical distribution function with averaging. For $N$ non-missing forecasts sorted
ascending as $x_1,\dots,x_N$, with $p = q/100$, $i = \lfloor Np \rfloor$ and $g = Np - i$,
\begin{equation}
\hat{y}_{t,q} =
\begin{cases}
\tfrac{1}{2}\left(x_i + x_{i+1}\right), & g = 0,\\[2pt]
x_{i+1}, & g > 0 .
\end{cases}
\end{equation}
The evaluation uses thirty years of historical temperature (1976--2005) held strictly prior
to the 2006--2011 test years, so that no scenario is constructed from weather observed within
the evaluation period.

\subsubsection{Quantile Regression Averaging on Sister Forecasts}
Approaching PLF from a forecast combination perspective, Liu et al.
\cite{liu2017probabilistic} apply quantile regression averaging (QRA) to a set of
\emph{sister point forecasts}. Sister forecasts are produced by a single model family whose
members differ only in the outcome of variable selection, so that their forecasts are highly
correlated and their disagreement is treated as the carrier of information about uncertainty.
The family is Tao's Vanilla benchmark augmented with a recency effect,
\begin{equation}
y_t = \text{Vanilla}(t) + \sum_{d=1}^{D} f\!\left(\tilde{T}_{t,d}\right)
                        + \sum_{j=1}^{J} f\!\left(T_{t-j}\right),
\end{equation}
where $f(\cdot)$ is the same month- and hour-interacted cubic polynomial applied to the base
temperature term, $T_{t-j}$ is the temperature $j$ hours earlier, and $\tilde{T}_{t,d}$ is the
mean temperature over the $d$-th preceding 24-hour block. Eight sisters are obtained as the
product of two estimation window lengths (two and three years) and four calibration schemes
for model selection, each scheme yielding its own recency order $(D,J)$.

In the second stage, the $m$ sisters with the lowest mean absolute error over a rolling
calibration window of $L$ days are selected, and a linear quantile regression
\cite{koenker1978regression} of the actual load on those sister forecasts is estimated
separately for each quantile level:
\begin{equation}
\hat{\beta}_q = \arg\min_{\beta}
  \sum_{t \in \mathcal{L}} \rho_q\!\left(y_t - \mathbf{x}_t^{\top}\beta\right),
\qquad
\rho_q(u) =
\begin{cases}
q\,u, & u \ge 0,\\
(q-1)\,u, & u < 0 ,
\end{cases}
\end{equation}
with $\mathbf{x}_t$ collecting an intercept and the selected sister forecasts. Estimation is
carried out in logarithmic space and repeated at every forecast origin, so the combination
weights are re-derived daily and a constant width in log space corresponds to a proportional
width in megawatts. The study evaluates QRA against three benchmarks constructed from the
same sister family---denoted Direct, Ind and BI---on the GEFCom2014 load track, reporting the
pinball loss and the Winkler score of the 50\% and 90\% prediction intervals.

\subsubsection{Combination of Probability Density Forecasts}
A third paradigm combines full predictive densities rather than point forecasts. Li et al.
\cite{li2020combining} assemble fifteen base learners in three families of five: Gaussian
process regressions, quantile regression neural networks, and quantile regression gradient
boosting machines. Diversity within each family is induced by a different random subsample
for the Gaussian processes, by two architectures crossed with different feature subsets for
the neural networks, and by different tree hyperparameters for the boosting machines. The
predictors are fifteen generic inputs---eleven lagged load observations together with hour of
day, day of week, month of year, and a weekend indicator---with no weather variables.

The Gaussian process family emits a Gaussian predictive density directly. The other two emit
deciles, which are converted to a Gaussian by noting that if the quantiles originate from
$\mathcal{N}(\mu,\sigma^2)$ then $y_q = \mu + p_q\sigma$ with $p_q = \Phi^{-1}(q)$, so that an
ordinary least squares fit of $y_q$ on $p_q$ recovers the mean as the intercept and the
standard deviation as the slope. The combined forecast is the Gaussian mixture
$F = \sum_i w_i F_i$, and the weights are chosen to minimize the continuous ranked probability
score (CRPS) \cite{gneiting2007strictly} over a calibration period.

The contribution that makes this tractable is that the CRPS of a Gaussian mixture is
available in closed form. Applying the identity of Baringhaus and Franz
\cite{baringhaus2004new},
\begin{equation}
\mathrm{CRPS}(F, y) = \mathbb{E}\left|Y - y\right|
                    - \tfrac{1}{2}\,\mathbb{E}\left|Y - Y'\right|,
\end{equation}
with $Y, Y'$ independent draws from $F$, and using the fact that a difference of Gaussian
mixtures is again a Gaussian mixture, the score becomes quadratic in the weights,
\begin{equation}
\mathrm{CRPS}(F, y) = \sum_i \sum_k \alpha_{i,k}\, w_i w_k + \sum_i \beta_i w_i ,
\end{equation}
where $\alpha_{i,k} = -\tfrac{1}{2}\mathbb{E}|Y_i - Y'_k|$ and
$\beta_i = \mathbb{E}|Y_i - y|$, both expectations of the absolute value of a normal variate
and therefore expressible in elementary functions. Determining the weights is consequently a
quadratic program on the simplex,
\begin{equation}
\min_{\mathbf{w}} \; \mathbf{w}^{\top}\mathbf{Q}\mathbf{w} + \mathbf{c}^{\top}\mathbf{w}
\quad \text{s.t.} \quad \mathbf{1}^{\top}\mathbf{w} = 1, \;\; \mathbf{w} \ge 0,
\end{equation}
rather than a simulation. Because the weights are constrained to be non-negative, a component
receiving zero weight is removed from the mixture, so model selection and weighting are
performed in a single solve. The diagonal entries of $\mathbf{Q}$ satisfy
$\alpha_{i,i} = -\sigma_i/\sqrt{\pi} < 0$; the authors address the resulting indefiniteness by
eliminating the equality constraint and report the reduced matrix to be positive definite in
all of the cases they examine. The data are partitioned into four consecutive periods for
base model training, base model selection, weight optimization, and testing, and the
CRPS-guided weights are compared against weights obtained by minimizing the MAPE of the
mixture mean and against a simple average. A related combination framework operating on
quantile rather than density forecasts is developed in \cite{wang2019combining}.

\subsubsection{Quantile Regression Neural Networks}
The fourth paradigm estimates the conditional quantile function directly. The quantile
regression neural network (QRNN) of Cannon \cite{cannon2011quantile} replaces the squared
error objective of a feedforward network with the pinball loss. Setting the width of the
output layer equal to the cardinality of the target quantile set $\mathcal{Q}$ and summing the
loss across it allows a single network to produce all quantiles in one forward pass:
\begin{equation}
\mathcal{L} = \frac{1}{|\mathcal{Q}|\,T}
  \sum_{q \in \mathcal{Q}} \sum_{t=1}^{T}
  \max\!\left\{\left(\tfrac{q}{100}-1\right)\left(y_t - \hat{y}_{t,q}\right),\;
               \tfrac{q}{100}\left(y_t - \hat{y}_{t,q}\right)\right\}.
\end{equation}

Zhang et al. \cite{zhang2019improved} propose an improved variant (iQRNN) in which the
training scheme of the canonical implementation is replaced by established deep learning
practice. Mini-batch training is used with a batch of $24 \times 7$ samples, chosen so that
each update spans a full weekly cycle; dropout and additive Gaussian noise layers of standard
deviation $0.3$ follow each of two hidden layers with parametric rectified linear activation;
training is terminated once the validation quantile score has not improved for $k$ epochs; and
the learning rate is multiplied by a decay factor $D_e$ whenever it has not improved for
$r < k$ epochs, subject to a floor. Feature selection is performed by recursive feature
elimination, in which random forests trained under the mean absolute error criterion rank
1760 candidate predictors---calendar and holiday indicators, load lags from one day to five
weeks, and temperature from the current time to five weeks earlier---and the least important
are pruned recursively until 100 remain, to which weekday and holiday indicators are added
manually, giving 103 inputs per sample.

The study evaluates the proposal on the GEFCom2014 load track \cite{hong2016gefcom} against
QRA \cite{liu2017probabilistic}, quantile regression forests in three implementations, and
gradient boosting regression trees in two, and against a baseline QRNN that retains the
parametric rectified linear activation and batch training but none of the four additions.
Each network configuration is trained one hundred times independently and reported as a mean
and standard deviation. The authors additionally report sensitivity analyses over network
width in the range 200--800 units per layer and over the Gaussian noise standard deviation in
the range $0$--$0.4$, and report training times of a few minutes for iQRNN against 85 hours
for the composite QRNN available in the reference \texttt{qrnn} R package.

\subsection{Deep Learning Approaches to Probabilistic Load Forecasting}
A substantial body of more recent work applies modern deep architectures to PLF. One line
addresses \emph{epistemic} uncertainty---arising from finite training data---separately from
the \emph{aleatoric} uncertainty inherent in the load process. Yang et al.
\cite{yang2020bayesian} cluster smart meter customers by their normalized weekly load
profiles, pool the profiles within each cluster, and train a multi-task Bayesian recurrent
network in which long short-term memory layers are shared across clusters while the output
layer is task-specific; predictive moments are obtained from repeated stochastic forward
passes under Monte Carlo dropout, with an additive term for inherent data noise estimated on a
separate validation set. Brusaferri et al. \cite{brusaferri2022probabilistic} replace the
deterministic output layer with a mixture density layer of spherical Gaussian kernels, whose
mixing coefficients, means and variances are all functions of the inputs, and place a
mean-field variational posterior over the network weights; deep ensembles are combined with
the variational approximation to cover several modes of that posterior, and the two
uncertainty components are separated through the differential entropy of the predictive
distribution and its mutual information with the weights.

A parallel line adapts attention-based sequence architectures. The temporal fusion
transformer \cite{lim2021temporal} couples variable selection networks, static covariate
encoders, and interpretable multi-head attention with a quantile output layer for
multi-horizon forecasting. Li et al. \cite{li2023probabilistic} adapt it to mid-term hourly
load forecasting by reconstructing the univariate hourly series into a panel of day-to-day
series indexed by hour of day, substituting gated recurrent units for the long short-term
memory layers, and augmenting the quantile loss with penalty terms for quantile crossing and
for prediction interval width. Faustine and Pereira \cite{faustine2022fpseq2q} parameterize
the quantile fractions themselves in addition to the quantile values, using one network to
propose the fractions through a cumulative softmax and a second, with a cosine embedding of
those fractions, to produce the corresponding values, trained jointly under a Huber quantile
loss together with a Wasserstein term.

\subsection{Evaluation of Probabilistic Forecasts}
Evaluation in this literature is based on strictly proper scoring rules
\cite{gneiting2007strictly}. The quantile score, the pinball loss averaged over the target
quantile set and the forecast horizon, is the metric adopted by GEFCom2014
\cite{hong2016gefcom} and is reported by all four of the methods above. It is commonly
accompanied by the Winkler interval score \cite{winkler1972decision}, which for a central
$(1-\alpha)\times 100\%$ prediction interval with bounds $L_t$ and $U_t$ and width
$\Delta_t = U_t - L_t$ is
\begin{equation}
\mathrm{WS}_t =
\begin{cases}
\Delta_t, & L_t \le y_t \le U_t,\\
\Delta_t + 2(L_t - y_t)/\alpha, & y_t < L_t,\\
\Delta_t + 2(y_t - U_t)/\alpha, & y_t > U_t,
\end{cases}
\end{equation}
so that width is charged at every step while the penalty for a miss is scaled by $1/\alpha$.
The CRPS \cite{gneiting2007strictly} summarizes the whole predictive distribution and is the
headline metric of the density combination literature; it coincides with twice the integral
of the pinball loss over the quantile level, which allows it to be evaluated from a quantile
function as well as from a density. Empirical coverage against the nominal level is widely
reported as a calibration diagnostic \cite{nowotarski2018recent}, and the ITFT study
\cite{li2023probabilistic} additionally defines a constraint score measuring the extent to
which adjacent predicted quantiles violate monotonicity.

\subsection{Summary}
Table~\ref{tab:baselines} summarizes the four baseline methods along the dimensions that
distinguish them: the source from which uncertainty is derived, the predictors entering the
model, whether parameters are re-estimated at each forecast origin, and the form of the
emitted forecast.

% TODO: 本节之后接你的研究问题/贡献。上表只做客观归纳，未作评价。

\begin{table}[!t]
\caption{Summary of the four probabilistic load forecasting baselines.}
\label{tab:baselines}
\centering
\small
\setlength{\tabcolsep}{4pt}
\begin{tabular}{@{}lllll@{}}
\toprule
Method & Uncertainty source & Predictors & Re-estimation & Output \\
\midrule
Scenario \cite{xie2018temperature} & Input weather ensemble & Calendar, temperature
  & Static  & Empirical percentiles \\
QRA \cite{liu2017probabilistic}    & Sister disagreement    & Calendar, temperature
  & Rolling & Quantiles \\
Density \cite{li2020combining}     & Component densities    & Load lags, calendar
  & Static  & Gaussian mixture \\
iQRNN \cite{zhang2019improved}     & Pinball objective      & Load lags, temperature, calendar
  & Static  & Quantiles \\
\bottomrule
\end{tabular}
\end{table}

% --- 参考文献渲染 ---
\bibliographystyle{IEEEtran}
\bibliography{\jobname}

\end{document}
